\subsection{Streaming Large Inputs}

Large input files should not always be loaded into memory as one complete object. A small text file can be handled with \texttt{read()}, and a small binary file can be handled with \texttt{read()} as well. A large log file, dataset export, image archive, or media file should usually be processed incrementally.

Streaming means that the program reads one manageable piece at a time, processes that piece, and then moves to the next piece. The whole file does not need to exist as one giant Python object in memory.

\subsubsection{Processing Files Line by Line Without Loading Entire Contents into Memory}

For text files, the most common streaming pattern is direct iteration over the file object. Each loop step receives one line as a \texttt{str} object. The program processes that line and then continues.

\begin{verbatim}
from pathlib import Path

path = Path("large_text_stream_demo.txt")

with open(path, "w", encoding="utf-8") as file_obj:
    file_obj.write("001 | Homer | D'oh!\n")
    file_obj.write("002 | Jerry | Serenity now!\n")
    file_obj.write("003 | Arthur | Don't panic.\n")
    file_obj.write("004 | Chuck | Counted to 42 twice.\n")

line_count = 0
char_count = 0

with open(path, "r", encoding="utf-8") as file_obj:
    print("File object used as a line stream:")
    print(f"file_obj       = {file_obj}")
    print(f"type(file_obj) = {type(file_obj)}")

    print()
    print("Selected names from dir(file_obj):")
    for name in ["__iter__", "read", "readline", "readlines", "close", "mode"]:
        print(f"{name:12} -> {name in dir(file_obj)}")

    print()
    print("Streaming lines:")

    for line in file_obj:
        line_count += 1
        char_count += len(line)

        clean_line = line.rstrip("\n")

        print(f"line {line_count:02d}: {clean_line}")
        print(f"    type(line) = {type(line)}")
        print(f"    len(line)  = {len(line)}")

print()
print("Summary:")
print(f"line_count = {line_count}")
print(f"char_count = {char_count}")
\end{verbatim}

The loop variable \texttt{line} refers to one \texttt{str} object at a time. The program does not build a list of all lines. It does not create one large string containing the entire file. It advances through the file sequentially.

This is the practical difference between streaming and materialization:

\begin{itemize}
    \item \texttt{file\_obj.read()} materializes the whole remaining file as one string;
    \item \texttt{file\_obj.readlines()} materializes all remaining lines as a list;
    \item \texttt{for line in file\_obj} processes one line per loop step.
\end{itemize}

Line streaming is useful for filters, counters, searches, log processors, and text transformations.

\begin{verbatim}
from pathlib import Path

path = Path("large_text_filter_demo.txt")

with open(path, "w", encoding="utf-8") as file_obj:
    file_obj.write("INFO  | Startup complete.\n")
    file_obj.write("WARN  | Disk almost full.\n")
    file_obj.write("INFO  | User asked for 42.\n")
    file_obj.write("ERROR | Nobody expects the Spanish Inquisition.\n")
    file_obj.write("INFO  | Shutdown complete.\n")

error_count = 0
warn_count = 0

with open(path, "r", encoding="utf-8") as file_obj:
    for line_number, line in enumerate(file_obj, start=1):
        clean_line = line.rstrip("\n")

        if clean_line.startswith("ERROR"):
            error_count += 1
            print(f"error at line {line_number:02d}: {clean_line}")

        if clean_line.startswith("WARN"):
            warn_count += 1
            print(f"warning at line {line_number:02d}: {clean_line}")

print()
print("Streaming summary:")
print(f"error_count = {error_count}")
print(f"warn_count  = {warn_count}")
\end{verbatim}

The program only keeps counters and the current line. That is the important memory property. The file may contain five lines or five million lines; the loop shape remains the same.

\subsubsection{Chunked Binary Processing for Large Media or Archive Files}

Binary files usually do not have lines. A binary stream is more naturally processed in fixed-size chunks. The program reads a block of bytes, processes that block, and repeats until \texttt{read(size)} returns \texttt{b""}.

\begin{verbatim}
from pathlib import Path

path = Path("large_binary_stream_demo.bin")

data = bytearray()

for number in range(1, 65):
    data.append(number)

with open(path, "wb") as file_obj:
    file_obj.write(data)

chunk_size = 10
chunk_count = 0
byte_count = 0

with open(path, "rb") as file_obj:
    print("Binary file object used as a byte stream:")
    print(f"file_obj       = {file_obj}")
    print(f"type(file_obj) = {type(file_obj)}")

    print()
    print("Selected names from dir(file_obj):")
    for name in ["read", "write", "seek", "tell", "close", "mode"]:
        print(f"{name:10} -> {name in dir(file_obj)}")

    print()
    print("Reading binary chunks:")

    while True:
        chunk = file_obj.read(chunk_size)

        if chunk == b"":
            break

        chunk_count += 1
        byte_count += len(chunk)

        print(f"chunk {chunk_count:02d}: {chunk}")
        print(f"    type(chunk) = {type(chunk)}")
        print(f"    len(chunk)  = {len(chunk)}")

print()
print("Binary streaming summary:")
print(f"chunk_count = {chunk_count}")
print(f"byte_count  = {byte_count}")
\end{verbatim}

The object returned by each \texttt{read(chunk\_size)} call is a \texttt{bytes} object. It is immutable and contains at most \texttt{chunk\_size} bytes. The final chunk may be smaller if the file ends before a full block is available.

Chunked processing is the normal shape for large binary transfers because it bounds memory use.

\begin{verbatim}
from pathlib import Path

src_path = Path("binary_source_demo.bin")
dst_path = Path("binary_copy_demo.bin")

with open(src_path, "wb") as file_obj:
    file_obj.write(b"Spam and eggs.\n")
    file_obj.write(b"D'oh!\n")
    file_obj.write(b"The answer is 42.\n")

chunk_size = 8
total_count = 0

with open(src_path, "rb") as src_obj:
    with open(dst_path, "wb") as dst_obj:
        while True:
            chunk = src_obj.read(chunk_size)

            if chunk == b"":
                break

            written_count = dst_obj.write(chunk)
            total_count += written_count

            print(f"copied {written_count:02d} bytes: {chunk!r}")

print()
print(f"total_count = {total_count}")

with open(dst_path, "rb") as file_obj:
    copied_data = file_obj.read()

print()
print("Copied data:")
print(copied_data)

print()
print("Object inspection:")
print(f"type(copied_data) = {type(copied_data)}")
print(f"len(copied_data)  = {len(copied_data)}")
\end{verbatim}

This example copies a binary file without reading the whole source file at once. The program repeatedly reads a small \texttt{bytes} chunk from the source and writes that same chunk to the destination.

A chunk can also be inspected or counted without being decoded as text.

\begin{verbatim}
from pathlib import Path

path = Path("binary_count_demo.bin")

with open(path, "wb") as file_obj:
    file_obj.write(b"AAAABBBBCCCC")
    file_obj.write(bytes([0, 1, 2, 3, 42, 255]))
    file_obj.write(b"AAA")

chunk_size = 5
target_byte = ord("A")
target_count = 0
total_bytes = 0

with open(path, "rb") as file_obj:
    while True:
        chunk = file_obj.read(chunk_size)

        if chunk == b"":
            break

        total_bytes += len(chunk)
        target_count += chunk.count(target_byte)

        print(f"chunk = {chunk!r}")
        print(f"    count of byte {target_byte} = {chunk.count(target_byte)}")

print()
print("Final counts:")
print(f"total_bytes  = {total_bytes}")
print(f"target_count = {target_count}")
\end{verbatim}

The variable \texttt{target\_byte} is an integer byte value. The \texttt{bytes.count()} method can count that value inside each chunk. No text decoding is involved.

The practical streaming rule is direct: text streams are usually processed line by line; binary streams are usually processed chunk by chunk. Both patterns avoid unnecessary full-file materialization and keep memory use tied to the current unit of work rather than the total file size.